\textbf{Dinic's Algorithm(Max Flow)}
\begin{lstlisting}
const int64_t inf = 1LL << 61;
struct Dinic {
  struct edge {
    int to, rev;
    int64_t flow, w;
  };
  int n, s, t;
  vector<int> d;
  vector<int> done;
  vector<vector<edge>> g;
  Dinic() {};
  Dinic(int _n) {
    n = _n + 10, g.resize(n);
  }
  void add_edge(int u, int v, int64_t w) {
    edge a = {v, (int)g[v].size(), 0, w};
    edge b = {u, (int)g[u].size(), 0, 0};
    g[u].push_back(a), g[v].push_back(b);
  }
  bool bfs() {
    d.assign(n, -1), d[s] = 0;
    queue<int> Q;
    Q.push(s);
    while (Q.size()) {
      int u = Q.front(); Q.pop();
      for (auto &e : g[u]) {
        int v = e.to;
        if (!~d[v] && e.flow < e.w) d[v] = d[u] + 1,
        Q.push(v);
      }
    }
    return d[t] != -1;
  }
  int64_t dfs(int u, int64_t flow) {
    if (u == t) return flow;
    for (int &i = done[u]; i < (int)g[u].size(); i++) {
      edge &e = g[u][i];
      if (e.w <= e.flow) continue;
      int v = e.to;
      if (d[v] == d[u] + 1) {
        int64_t nw = dfs(v, min(flow, e.w - e.flow));
        if (nw > 0) {
          e.flow += nw;
          g[v][e.rev].flow -= nw;
          return nw;
        }
      }
    }
    return 0;
  }
  int64_t max_flow(int _s, int _t) {
    s = _s, t = _t;
    int64_t flow = 0;
    while (bfs()) {
      done.assign(n, 0);
      while (int64_t nw = dfs(s, inf)) flow += nw;
    }
    return flow;
  }
};
\end{lstlisting}
